\chapter*{Preface}
\addcontentsline{toc}{chapter}{Preface}

A compiler changes the form of a program without losing its meaning.
That short promise leads to many questions.
Which parts of an execution matter?
What can we learn without running a program?
When does a shorter instruction sequence do more work?
How do we find a mistake that appears only after several correct-looking changes?

This book approaches those questions by building things small enough to inspect.
We start with programs whose executions fit on a page.
We give names to the facts their translations need.
Then we develop representations, analyses, transformations, and runtime services that let the same reasoning reach larger programs.
The aim is not to memorize a catalog of passes.
It is to understand what each pass promises and how an implementation keeps that promise.

The reader should know basic data structures, elementary discrete mathematics, and one systems language.
Executable examples use C11 unless an example states a narrower requirement.
Algorithms use pseudocode so that their contracts remain separate from the details of C.
Mathematical notation is introduced where it is needed, next to an example.

The accompanying pamphlet, \emph{PebbleIR: Design, Development and Dogfooding}, develops a small executable intermediate representation in sixteen chapters.
Its independent master file is \code{book/pamphlet.ltx}; it shares this book's typography and can be built as a separate companion PDF.
The C11 implementation connects construction, verification, interpretation, rewriting, and lowering in one tested path.

The reference books provide background, not text to reproduce.
Their ideas are discussed in original prose, with new examples and citations to the relevant works.
Claims about a particular language follow its specification rather than a convenient analogy with algebra.
Claims about performance require measurement rather than confidence.

\begin{ccnote}[Manuscript status]
All chapters and appendices are present as drafts under technical review.
Chapter lengths were chosen for their content rather than a uniform word target.
Runnable examples currently cover the opening laboratory and selected review checks.
Later implementation exercises are assignments, not claims of tested implementations.
Examples, measurements, and primary-source citations still need further verification.
\end{ccnote}
